\section{Related Work and Background}
\vspace{-1mm}
\subsection{Post-Training Compression Method}
\vspace{-1mm}
Under inference constraints on accuracy, latency, and memory, \emph{post-training} compression has become a practical approach to obtaining model instances by directly producing executable compressed models from LLM bases~\cite{han2016deepcompression,frankle2019lottery,wan2023efficientllm_survey}.
Typically, retraining-free unstructured pruning methods, such as SparseGPT~\cite{frantar2023sparsegpt} and Wanda~\cite{sun2023wanda}, estimate importance from weights or activations while still relying on immense computational burdens.
Meanwhile, recent structured pruning emphasizes executable budgets (e.g., layers, blocks, or heads), yielding more performance preserved and memory benefits~\cite{llmpruner2023,ashkboos2024slicegpt,yao2024gpruner,he2025olica}.
However, these approaches often rely on \emph{global} or \emph{average} importance/reconstruction objectives, making pruned models sensitive to diverse scenarios ~\cite{michel2019sixteen,lele2025rethinking_tfsp,yuan2023revisitingoodnlp,yang2024advoodrobustness}.
In contrast, we move to \emph{scenario-specific} pathway calculation, using embedding-space domain axes as a coordinate system to identify and preserve axis-aligned pathways, yielding an executable pruned instance for specialized reuse.
\vspace{-0.5mm}

% \subsection{Computational Pathways and Structured Pruning}
% Our structure-preserving principle aligns with evidence that effective Transformer computation concentrates in a small and stable subset of components, motivating pathway-level calculation rather than relying on static weights~\cite{michel2019sixteen,voita2019heads}.
% Related mechanistic analysis and interpretability studies further suggest that ``behaviors'' can be associated with localized internal ``mechanisms'', and that sparsity can make such structure more traceable~\cite{elhage2021framework,elhage2022superposition,openai2025sparsecircuitsblog}.
% %Moreover, findings on feature entanglement (e.g., superposition/polysemanticity) highlight that different behaviors may overlap in shared components, which can make naive static pruning brittle under shift \cite{elhage2022superposition}.
% These results suggest that specialization-oriented pruning should not only remove global redundancy but also preserve the key \emph{reasoning pathways} supporting target ``behavior''.
% While prior work leverages these insights only for analysis or generic compression, our method turns them into an operational criterion via \emph{representation--parameter coupling}. 
%Specifically, while embedding-space axes provide a stable reference system, our method conservatively preserves pathway sets aligned with target axes.
\subsection{Structural Sparsity and Reasoning Pathways}
\vspace{-1mm}
Our approach builds on the observation that Transformer computation is often concentrated in sparse, task-specific subsets of components rather than utilizing all static weights~\cite{michel2019sixteen,voita2019heads}. Mechanistic interpretability studies further suggest that specific model ``behaviors'' are driven by localized internal ``mechanisms'', and that enforcing sparsity can make these structures more traceable~\cite{elhage2021framework,openai2025sparsecircuitsblog}. While prior work on structured pruning focuses on global redundancy removal or generic compression~\cite{frantar2023sparsegpt}, the challenge remains in how to dynamically activate the specific ``reasoning pathways'' that support a target behavior during inference. Our method addresses this by treating these pathways not as static parameters, but as functional units that can be selectively engaged.
In parallel, our method shares the input-dependent sparsity principle with mixture-of-experts (MoE), but targets dense-checkpoint specialization~\citep{shazeer2017outrageously,lepikhin2020gshard,fedus2022switch,jiang2024mixtral}. Unlike MoE token-level routing over FFN experts, we compiles a reusable scenario-level head mask from a fixed dense model.
\vspace{-0.5mm}

\subsection{Representation Subspaces and Probing}
\vspace{-1mm}
Another pillar of our work is the internal structure of the LLM's embedding space. Research in representation probing has long demonstrated that task-relevant information is often encoded in linear subspaces or axes~\cite{alain2016linearprobes,belinkov2022probing}. Recent advances in sparse feature modeling (e.g., SAEs) further reveal that these internal activations form a partially separable manifold that can be recovered to identify specific semantic domains~\cite{gao2024sae_scaling,bricken2023dictionary}. While recent ``probe-driven pruning''~\cite{le2025probe} uses probing signals to guide model reduction, it often lacks a formal grounding in representation alignment, leading to decisions that are difficult to generalize across diverse scenarios. By contrast, we leverage the stability of these latent subspaces as "coordinates" to provide the necessary signaling for our pruning mechanism.

% Another key evidence of our approach is that representation subspaces/axes in the embedding space can serve as actionable coordinates, supporting lightweight probing signals to realize precise scenario-specific computing.
% Linear probes have long been used to assess whether task-relevant information is linearly readable from representations~\cite{alain2016linearprobes,belinkov2022probing}. Recent advances in representation decomposition and sparse feature modeling provide further evidence that internal activations exhibit a partially separable structure that can be recovered and evaluated~\cite{gao2024sae_scaling,bricken2023dictionary,henighan2024scalingmonosemanticity}.
% The most closely related work is probe-driven pruning, exemplified by Probe Pruning~\cite{le2025probe}, which shows that probing signals can guide pruning decisions.
% However, such methods do not explicitly ground probe outputs in \emph{representation alignment}, making it harder to interpret and reuse pruning decisions across diverse scenarios.
% By contrast, our method anchors probes to a quasi-orthogonal axis basis in embedding space and uses stable and reliable probe-derived signals to robustly compile \emph{pathway} within inference time.